\documentclass{zapiski}
\originfo{1}{}{}{2026}
\setcounter{page}{1}
\date{}

\usepackage[utf8]{inputenc}
\usepackage[T2A]{fontenc}
\usepackage[russian, english]{babel}

% Fix the incorrectly decoded Russian masthead from zapiski.cls
\renewcommand{\publnamef}{%
  \hbox{\foreignlanguage{russian}{Записки научных}}}
\renewcommand{\publnames}{%
  \hbox{\foreignlanguage{russian}{семинаров ПОМИ}}}
\renewcommand{\russianvolinfo}{%
  \hbox{\foreignlanguage{russian}{Том~\origvolume, \origyear~г.}}}

\usepackage{times}
\usepackage{url}
\usepackage{latexsym}
\usepackage{multirow}
\usepackage{graphicx}
\usepackage{amssymb}
\usepackage{amsmath}
\usepackage{hyperref}
\usepackage{booktabs}
\usepackage{tabularx}
\usepackage{caption}
\usepackage{cite}
\usepackage{paralist}
\usepackage{subcaption}
\usepackage{xcolor,colortbl}
\usepackage{array}
\usepackage{setspace}
\usepackage{placeins}
\usepackage{tikz}
\usepackage{pgfplots}
\pgfplotsset{compat=1.17}

\newcolumntype{H}{>{\setbox0=\hbox\bgroup}c<{\egroup}@{}}
\renewcommand{\UrlFont}{\ttfamily\small}
\newcommand{\cm}[1]{\textcolor{red}{#1}}
\newcommand{\bm}[1]{\textcolor{blue}{#1}}

\usepackage{float}
\restylefloat{table}

% --- notation shorthands -------------------------------------------------
\newcommand{\pmap}[2]{\pi^{\theta}_{#1,#2}}
\newcommand{\pAR}{p_{\mathrm{AR}}}
\newcommand{\sg}[1]{\mathrm{sg}\!\left[#1\right]}

\english

\begin{document}

\title[FlowDraft]{FlowDraft: A Categorical Flow-Map Drafter for Lossless Parallel Decoding}

\author[Sergaev et al.]{Ya.~Sergaev, V.~Tekaev, N.~Nikonov}
\address[Ya.~Sergaev]{\cm{affiliation}}
\email{\cm{email}}

% \cm{The call requires a named per-member Contribution section, so this
% submission is not anonymous. Confirm on OpenReview before uploading; if the
% track turns out to be double-blind, the Contribution section stays and the
% author block is replaced, not the other way round --- the section is a
% stated requirement.}

\begin{abstract}
Speculative decoding removes the sequential bottleneck of autoregressive (AR)
language models without changing their output distribution: a cheap drafter
proposes a block of tokens in parallel and the frozen AR model verifies it.
Throughput is then governed by a single quantity --- the number of drafted
tokens accepted per cycle. We replace the single-step masked-diffusion drafter
of the Orthrus framework with a drafter parameterised as a categorical flow
map, which produces a correlated joint proposal over the block at the same
number of forward passes. On a frozen Qwen3-1.7B backbone our drafter reaches
$2.07$ accepted tokens per cycle and $1.44$ tokens per forward pass, against
$2.00$ and $1.39$ for the reproduced Orthrus baseline, with verification
guaranteeing bit-identical output in both cases. Two findings concern where
that performance comes from. Training the drafter in the exact geometry it
decodes in --- clean prefix cache, clean in-block anchor, noised block --- is
what makes it work at all: under otherwise identical settings, a drafter
trained on fully noised sequences accepts nothing, while the block-wise
drafter accepts $0.412$ tokens per cycle. And along the training trajectory,
$92\%$ of the improvement is realised before the self-distillation terms reach
their full weight, which bounds what may presently be attributed to the
flow-map formulation itself.
\end{abstract}

\keywords{speculative decoding \and flow maps \and discrete diffusion \and
efficient inference \and language models.}

\maketitle

% =========================================================================
\section{Introduction}
\label{sec:intro}

Autoregressive decoding emits one token per forward pass. At inference the
cost of that pass is dominated by moving weights from memory rather than by
arithmetic, so a language model generating $L$ tokens pays $L$ memory-bound
passes regardless of how much parallel compute is available. Diffusion
language models break the dependency by emitting a whole block at once, but
they factorise the block into conditionally independent positions and the
resulting text drifts away from the AR distribution.

Speculative verification reconciles the two. A drafter proposes a block, the
frozen AR model scores it in a single pass, and only the longest prefix the AR
model would itself have produced is kept. Output is then \emph{lossless} by
construction: identical to what the AR model alone would have emitted. Quality
is no longer a design variable, and the only quantity that matters is the
\emph{acceptance length} --- how many drafted tokens survive verification per
cycle.

The Orthrus framework~\cite{nguyen2026orthrus} realises this with a frozen AR
backbone augmented by a lightweight trainable diffusion view that shares the
same KV cache. Its drafter is a single-step masked diffusion model, so the
positions of a block are proposed independently of one another; a block whose
tokens are individually plausible but jointly incoherent is rejected after its
first mismatch. Raising acceptance therefore requires a better \emph{joint}
proposal at the same pass count, not more passes --- an extra refinement pass
costs exactly what it saves.

Categorical flow maps~\cite{roos2026cfm} learn an integrated, correlated
endpoint distribution on the probability simplex and can be sampled in one or
few jumps. This paper uses such a map as the drafter inside the Orthrus
verification scaffold. Our contributions are:

\begin{compactenum}
\item a flow-map drafter for lossless parallel decoding, trained in the exact
geometry it decodes in --- a clean AR prefix in the KV cache, a clean in-block
anchor, and a noised block (Section~\ref{sec:model});
\item a dual distillation objective combining AR-verifier alignment with the
flow-map endpoint and consistency terms, together with the observation that
placing the verifier term on the diagonal of the map conflicts with the
categorical endpoint objective, while placing it on the jump the decode loop
executes does not (Section~\ref{sec:objective});
\item a controlled ablation of that training geometry --- the single change
from full-sequence to block-wise training moves acceptance from zero to
$0.412$ at otherwise identical settings (Section~\ref{sec:ablation-geometry});
\item an attribution analysis over the training trajectory showing which term
of the objective actually drives acceptance, and an explicit statement of which
claims the present evidence does and does not support
(Section~\ref{sec:attribution}).
\end{compactenum}

% =========================================================================
\section{Related Work}
\label{sec:related}

\paragraph{Speculative decoding.} \cm{TODO: Leviathan et al., Chen et al.,
Medusa, EAGLE --- draft-then-verify, lossless guarantee, why acceptance length
is the figure of merit.}

\paragraph{Diffusion language models.} \cm{TODO: masked diffusion LMs, the
conditional-independence limitation.}

\paragraph{Orthrus.} Nguyen et al.~\cite{nguyen2026orthrus} augment a frozen
LLM with a trainable parallel diffusion view sharing the AR key--value cache;
an exact consensus mechanism between the two views makes generation lossless,
with reported speedups up to $7.8\times$ and $O(1)$ cache overhead. Their
drafter is trained with a single objective, the KL divergence from the frozen
AR distribution to the diffusion view's output. \cm{TODO: verify the $7.8\times$
figure and the training-token count against the paper before submission.}

\paragraph{Categorical flow maps.} Roos et al.~\cite{roos2026cfm} define a
two-time map $\pmap{s}{t}$ towards the simplex that transports probability mass
onto a predicted endpoint, trained with an endpoint likelihood on the diagonal
and an endpoint-consistency self-distillation objective off it. The
construction targets few-step and single-step generation. Subsequent
work~\cite{davis2026scaling} scales it to $1.7$B parameters and $2.1$T tokens,
reporting high-quality text in four inference steps. In both cases distillation
is \emph{self}-distillation: the target is the model's own diagonal under a
stop-gradient, with no external teacher. Our setting differs in that a frozen
AR verifier is present by construction, which makes an external teacher
available and changes the design space of the objective.

% =========================================================================
\section{Model Description}
\label{sec:model}

\subsection{Setup}

A frozen AR backbone $\pAR$ provides both the context representation and the
verifier. A lightweight trainable diffusion head (DF) shares the backbone's
key--value cache and adds a flow-time embedding. Let $V$ be the vocabulary and
$\Delta^{V-1}$ the probability simplex. Clean tokens are embedded as simplex
vertices $x_1 \in \{e_i\}$, the prior is $x_0 \sim \mathrm{Dirichlet}(\mathbf{1})$,
and the linear path is
\[
x_s = (1-s)\,x_0 + s\,x_1 , \qquad s \in [0,1].
\]
The drafter is a two-time flow map $\pmap{s}{t}(x_s) \in \Delta^{V-1}$: an
endpoint prediction issued at noise level $s$ for target level $t$. Transport
is the convex step
\[
X_{s,t}(x_s) = x_s + \gamma_{s,t}\bigl(\pmap{s}{t}(x_s) - x_s\bigr),
\qquad \gamma_{s,t} = \frac{t-s}{1-s}.
\]
Decoding with a schedule $0 = t_0 < \dots < t_n = 1$ applies $n$ such steps; a
one-jump schedule calls $\pmap{0}{1}(x_0)$ alone.

\subsection{Inference geometry at training time}
\label{sec:geometry}

A drafter trained on fully noised sequences without a cache never sees the
configuration it decodes in. We therefore train block-wise, so that every
optimisation step reproduces one decode cycle: a single no-grad AR pass builds
the clean prefix cache and the teacher distributions; position $p$ carries the
previous cycle's correction token as a \emph{clean in-block anchor} whose
key/value pairs are not yet in the cache; positions $p+1,\dots,p+K-1$ are
noised and drafted. Every DF forward of the loss runs with that cache and that
anchor. Several such blocks are packed into one pass and isolated by the
dual-pass attention mask, so that a block attends only to the AR prefix
preceding its own anchor and to itself bidirectionally.

\begin{figure}[!tbh]
    \centering
    \cm{TODO: figure --- decode cycle and the matching training block.
    Panels: (a) prefill and cache, (b) block draft over $n$ jumps,
    (c) verification and accepted prefix, (d) the same geometry at train time.}
    \caption{The drafter is trained in the geometry it decodes in.}
    \label{fig:model}
\end{figure}

\subsection{Verification}

Given a drafted block, one AR forward yields the distributions the verifier
would have produced at each position. Under greedy decoding a drafted token is
accepted iff it equals the verifier's argmax conditioned on the accepted
prefix; the first mismatch invalidates everything after it, and the verifier's
own token is emitted in its place. Sampling uses the standard coupled scheme
with shared Gumbel noise. In both cases output is bit-identical to the AR model
alone, which we assert in evaluation.

% =========================================================================
\section{Training Objective}
\label{sec:objective}

The objective combines a verifier term with the flow-map terms:
\begin{equation}
\label{eq:loss}
\mathcal{L} =
  w_{\mathrm{ver}}\,\mathcal{L}_{\mathrm{ver}}
+ w_{\mathrm{end}}\,\mathcal{L}_{\mathrm{end}}
+ \lambda\,\bigl(4\,\mathcal{L}_{\mathrm{EC}} + 2\,\mathcal{L}_{\mathrm{TD}}\bigr).
\end{equation}

\paragraph{Verifier alignment.}
$\mathcal{L}_{\mathrm{ver}} = \mathrm{KL}\bigl(\sg{\pAR} \,\|\, \pmap{0}{1}(x_0)\bigr)$,
evaluated on the exact one-jump map the decode loop executes. The KL input is
the drafter's output at $(s,t) = (0,1)$; the KL target is the frozen AR path's
distribution for the same block positions.

\paragraph{Endpoint likelihood.}
$\mathcal{L}_{\mathrm{end}} = \mathrm{CE}\bigl(x_1, \pmap{t}{t}(x_t)\bigr)$, the
categorical endpoint likelihood on the diagonal. This is the anchor the
flow-map construction requires: it makes $\pmap{t}{t}$ the denoiser associated
with the sampled interpolation.

\paragraph{Endpoint consistency and temporal drift.}
$\mathcal{L}_{\mathrm{EC}} = \mathrm{CE}\bigl(\sg{\pmap{t}{t}(X_{s,t}(x_s))},\,
\pmap{s}{t}(x_s)\bigr)$ requires the jump to agree with the stop-gradient
diagonal evaluated at the jump's own landing point, and
$\mathcal{L}_{\mathrm{TD}} = \gamma_{s,t}^2 \lVert \partial_t \pmap{s}{t}(x_s) \rVert^2$
keeps the family smooth in $t$. Together these are the self-distillation half:
they are what propagates the diagonal to the jumps without an external teacher.

\paragraph{Where the verifier term attaches.}
An earlier version of this objective placed the verifier KL on the diagonal,
$\mathrm{KL}(\sg{\pAR} \,\|\, \pmap{t}{t}(x_t))$, on the grounds that the
teacher should anchor the local denoiser and the jumps should inherit it
through consistency. This does not train: on a given trajectory the diagonal
then receives two incompatible targets, the soft teacher distribution and the
sampled categorical endpoint $x_1$. Training collapses --- decode acceptance
falls to zero while agreement with the teacher continues to improve, the
signature of a diagonal that learns while nothing propagates to the jumps.
Attaching the verifier term to $\pmap{0}{1}$ instead leaves one target per
object and trains stably. \cm{TODO: include the failed-run curve as a figure;
data is logged.}

\paragraph{Time sampling.} Pairs are drawn as $t \sim U[0,1]$, $s \sim U[0,t]$,
giving density $\propto 1/t$ on the triangle. We note for later reference that
the corner $(s,t) = (0,1)$ --- the operating point of a one-jump schedule --- is
a measure-zero event under this scheme, so $\mathcal{L}_{\mathrm{ver}}$ is the
only term supervising it directly.

% =========================================================================
\section{Dataset}
\label{sec:data}

We train on the Nemotron post-training corpus\footnote{\url{https://huggingface.co/datasets/nvidia/Nemotron-Post-Training-Dataset-v2}},
using the \texttt{chat}, \texttt{math} and \texttt{code} splits. Documents are
packed to sequences of $2048$ tokens; drafted spans are constrained not to
cross a document boundary, so a block never continues past an end-of-sequence
token into an unrelated example. \cm{TODO: exact train\_size actually consumed,
token count, validation split size and selection.}

Evaluation prompts are drawn from GSM8K, MATH-500, HumanEval, MBPP, AIME24 and
AIME25. \cm{TODO: number of prompts per benchmark, generation length, decoding
temperature; confirm which of these were actually evaluated --- current logs
contain validation-set decode metrics only.}

% =========================================================================
\section{Experiments}
\label{sec:experiments}

\subsection{Metrics}

\paragraph{Acceptance length} is the number of drafted tokens accepted per
decode cycle, the quantity the whole system optimises.

\paragraph{Tokens per forward pass (TPF)} is total emitted tokens divided by
forward passes. A cycle costs $n{+}1$ passes for an $n$-jump schedule, so with
$n = 1$ zero acceptance yields $\mathrm{TPF} = 0.5$ --- the analytic floor,
half the speed of plain AR decoding, which emits one token per pass
($\mathrm{TPF} \approx 1$). TPF is the throughput figure of merit because it is
hardware-independent; wall-clock speedup is reported as a diagnostic only.

\paragraph{Teacher agreement} is the fraction of block positions where the
drafter's argmax matches the verifier's, measured off-policy on validation
sequences. It correlates with but does not determine acceptance, which is
measured on-policy during real generation.

\subsection{Experiment setup}

The backbone is Qwen3-1.7B, frozen; only the DF projections, the flow-time
embedding and the associated norms are trained. Block size is $K = 8$ (one
clean anchor and seven drafted positions). Optimisation uses AdamW with a
cosine schedule and $5\%$ warmup. \cm{TODO: hardware, wall-clock per run,
precision, seed count --- all runs to date use a single seed (42), which must
be stated or fixed.}

\subsection{Baselines}

The baseline is our reproduction of Orthrus: the same frozen backbone and
verification loop, with the masked-diffusion drafter and its single KL
objective. \cm{TODO: report the validation of this reproduction against the
authors' released weights --- the harness exists but its results are not yet
logged. Without it the baseline column is our implementation, not the
published system, and must be labelled as such.}

% =========================================================================
\section{Results}
\label{sec:results}

\subsection{Main comparison}

\begin{table}[!tbh]
\centering
\begin{tabular}{lrrrr}
\toprule
Drafter & Steps & Acceptance & TPF & Teacher agr. \\
\midrule
AR decoding (reference)      & ---    & ---   & $\approx 1.00$ & --- \\
Orthrus (masked diffusion)   & 9\,506 & 2.002 & 1.387 & 0.552 \\
FlowDraft (flow map)         & 11\,037& \textbf{2.070} & \textbf{1.441} & 0.536 \\
\bottomrule
\end{tabular}
\caption{Best run of each drafter, Qwen3-1.7B backbone, $K=8$, one-jump
schedule. Output is verified bit-identical to AR decoding in both rows.}
\label{tab:main}
\end{table}

The flow-map drafter improves acceptance by $3.4\%$ and TPF by $3.9\%$ over the
reproduced baseline. Teacher agreement is slightly \emph{lower}, which is
consistent with acceptance being an on-policy quantity: what matters is
agreement along the trajectory the drafter itself generates, not on
teacher-forced validation sequences.

\subsection{Full experimental record}

Table~\ref{tab:allruns} lists every training run behind the comparison.

\begin{table}[!tbh]
\centering\small
\begin{tabular}{llrrrrrrrr}
\toprule
Drafter & Date & $K$ & anch. & lr & $w_{\mathrm{ver}}$ & $\lambda$ & Steps & Acc. & TPF \\
\midrule
\multicolumn{10}{l}{\emph{Orthrus baseline, $K=32$}}\\
Orthrus & 07-13 & 32 & 256 & 2e-4 & --- & 1.0 & 7\,422  & 1.917 & 1.280 \\
Orthrus & 07-14 & 32 & 256 & 2e-4 & --- & 1.0 & 9\,380  & 1.500 & 1.103 \\
Orthrus & 07-14 & 32 & 256 & 2e-4 & --- & 1.0 & 6\,030  & 1.692 & 1.185 \\
Orthrus & 07-14 & 32 & 256 & 2e-4 & --- & 1.0 & 17\,828 & 1.698 & 1.261 \\
Orthrus & 07-15 & 32 & 256 & 2e-4 & --- & 1.0 & 8\,487  & 1.882 & 1.249 \\
Orthrus & 07-16 & 32 & 256 & 2e-4 & --- & 1.0 & 18\,159 & 1.534 & 1.186 \\
\midrule
\multicolumn{10}{l}{\emph{Orthrus baseline, $K=8$ --- identical configuration, three runs}}\\
Orthrus & 07-17 & 8 & 256 & 2e-4 & --- & 1.0 & 9\,506 & 1.592 & 1.205 \\
Orthrus & 07-18 & 8 & 256 & 2e-4 & --- & 1.0 & 9\,506 & 1.493 & 1.156 \\
Orthrus & 07-20 & 8 & 256 & 2e-4 & --- & 1.0 & 9\,506 & 2.002 & 1.387 \\
\midrule
\multicolumn{10}{l}{\emph{FlowDraft without verifier alignment}}\\
full-seq    & 07-20 & 8  & 1 & 1e-4 & 0 & 1.0 & 209 & 0.000 & 0.508 \\
full-seq    & 07-20 & 8  & 1 & 1e-4 & 0 & 1.0 & 209 & 0.033 & 0.525 \\
block-wise  & 07-20 & 8  & 1 & 1e-4 & 0 & 1.0 & 588 & 0.182 & 0.599 \\
block-wise  & 07-20 & 8  & 1 & 1e-4 & 0 & 1.0 & 209 & 0.412 & 0.713 \\
block-wise  & 07-20 & 16 & 1 & 1e-4 & 0 & 1.0 & 209 & 0.348 & 0.681 \\
block-wise  & 07-23 & 8  & 1 & 2e-4 & 0 & 1.0 & 789 & 0.022 & 0.519 \\
\midrule
\multicolumn{10}{l}{\emph{FlowDraft with verifier alignment}}\\
block-wise & 07-26 & 8 & 64 & 1e-4 & 1.0 & 0.25 & 5\,414  & 2.008 & 1.409 \\
block-wise & 07-27 & 8 & 32 & 1e-4 & 1.0 & 0.10 & 11\,037 & \textbf{2.070} & \textbf{1.441} \\
\bottomrule
\end{tabular}
\caption{Every logged training run. ``anch.'' is
\texttt{anchors\_per\_sequence}, the number of independently anchored blocks
per packed 2048-token sequence. TPF $=0.5$ is the analytic floor at zero
acceptance for a one-jump schedule.}
\label{tab:allruns}
\end{table}

Three points are visible in the table and must be addressed rather than
omitted.

\paragraph{Run-to-run spread exceeds the reported effect.} The three $K=8$
baseline runs share an identical configuration and step count and yield TPF
$1.205$, $1.156$ and $1.387$ --- a spread of $0.231$, or $\pm 9\%$ about their
mean of $1.249$. The claimed FlowDraft improvement over the \emph{best} of them
is $0.054$. A single run cannot be separated from baseline noise at this
magnitude.

\cm{This is the primary threat to the paper. Its cause is almost certainly the
measurement: decode metrics are computed over
\texttt{val\_decode\_prompts}~$=16$ prompts. Sixteen prompts cannot resolve a
$4\%$ difference in acceptance. Required: re-evaluate the existing checkpoints
on the held-out benchmarks with a sample size large enough to report a
confidence interval, and report mean $\pm$ std over seeds rather than the best
run. This costs evaluation compute only --- no retraining --- and is the single
highest-value item remaining.}

\paragraph{Verifier alignment is necessary in the present recipe.} No run
without $\mathcal{L}_{\mathrm{ver}}$ exceeds TPF $0.713$; several sit at the
floor of $0.5$. This is suggestive but not conclusive: those runs are $209$ to
$789$ steps against $11\,037$ for the main run, so they establish only that the
flow-map terms alone do not train quickly, not that they cannot train.

\paragraph{The main comparison is not compute-matched.} The baseline uses
\texttt{anchors\_per\_sequence}~$=256$ and lr~$2\!\times\!10^{-4}$; FlowDraft
uses $32$ and $1\!\times\!10^{-4}$, and its run terminated on a crash rather
than completing. The eightfold difference in blocks seen per optimiser step
favours the baseline, so the comparison is conservative in that respect and
optimistic in others; either way it is not controlled.

\subsection{Ablation: training in the inference geometry}
\label{sec:ablation-geometry}

The two runs in Table~\ref{tab:geometry} differ in exactly one field of the
configuration --- the drafter variant. Block size, anchors per sequence,
learning rate, all objective weights, the ramp schedule, time sampling, batch
size, sequence length, data splits, seed, backbone and step count are
identical, and both completed normally.

\begin{table}[!tbh]
\centering
\begin{tabular}{lrr}
\toprule
Training geometry & Acceptance & TPF \\
\midrule
Full sequence, no cache, no anchor   & 0.000 & 0.508 \\
Block-wise (inference geometry)      & \textbf{0.412} & \textbf{0.713} \\
\bottomrule
\end{tabular}
\caption{Effect of training in the decode configuration. Both runs: $K=8$,
one anchor per sequence, lr $1\!\times\!10^{-4}$,
$w_{\mathrm{end}} = \lambda = 1$, $w_{\mathrm{ver}} = 0$, seed 42, 209 steps.
The only difference is whether the DF forward carries the clean prefix cache
and the clean in-block anchor. TPF $0.508$ is the analytic floor.}
\label{tab:geometry}
\end{table}

Trained on fully noised sequences without a cache, the drafter accepts nothing
at all: every drafted block is rejected at its first position and TPF sits at
the floor. Reproducing the decode configuration at training time --- clean AR
prefix in the cache, clean in-block anchor, noised block --- moves acceptance
off zero without any change to the objective. Varying the block width in the
same setting gives $0.412$ at $K=8$ against $0.348$ at $K=16$, consistent with
longer blocks being harder to propose jointly.

\cm{Caveat that must stay until it is fixed: these runs are 209 steps and their
decode metrics use \texttt{val\_decode\_prompts}~$=2$. The qualitative signal
(zero versus non-zero acceptance) is clear, the point estimates are not. Rerun
this pair at the budget of the main experiment and with a proper evaluation
sample --- it is cheap, it is the only controlled comparison in the paper, and
it is the one a reviewer will trust most.}

\paragraph{What remains confounded.} The verifier alignment term is defined on
the one-jump map evaluated with the prefix cache and the in-block anchor, so it
exists only in the block-wise variant; every run that uses it is therefore also
a block-wise run. The final configuration consequently cannot be decomposed
into a geometry contribution and an objective contribution --- the two are
measured only together. Table~\ref{tab:geometry} isolates the geometry at
$w_{\mathrm{ver}} = 0$, which is the part that can be separated; separating the
remainder would require a full-sequence analogue of the verifier term, which is
a different object and is not implemented.

\subsection{What drives acceptance}
\label{sec:attribution}

The weight $\lambda$ on the self-distillation terms is ramped linearly from $0$
over the first phase of training, reaching its full value at step $3\,427$ of
the main run. This gives a natural staging of the objective, summarised in
Table~\ref{tab:phases}.

\begin{table}[!tbh]
\centering
\begin{tabular}{lrrrr}
\toprule
Phase & Steps & $\lambda$ & TPF & $\Delta$TPF \\
\midrule
Self-distillation effectively off & $0 \to 250$        & $0.001$--$0.007$ & $0.579 \to 1.095$ & $+0.516$ \\
$\lambda$ ramping                 & $250 \to 3\,427$   & $0.007 \to 0.1$  & $1.095 \to 1.372$ & $+0.277$ \\
$\lambda$ at full weight          & $3\,427 \to 11\,037$ & $0.1$          & $1.372 \to 1.441$ & $+0.069$ \\
\bottomrule
\end{tabular}
\caption{Training phases of the main FlowDraft run. Of the total TPF gain of
$0.862$, $60\%$ is realised in the first $250$ steps and $92\%$ before the
self-distillation weight reaches its full value.}
\label{tab:phases}
\end{table}

Two observations follow.

First, the endpoint term is not inert. It starts at $19.63$ and falls to
$0.011$, a drop of three orders of magnitude; at initialisation it accounts for
$31\%$ of the objective. A term driven to zero has been satisfied, not ignored.

Second, the self-distillation terms cannot be credited with the result on this
evidence. They reach full weight only after $92\%$ of the gain is banked, and
over the remaining $69\%$ of training they accompany a further $8\%$ of it. A
second run supports the same reading: reducing $w_{\mathrm{end}}$ from $0.5$ to
$0.25$ and $\lambda$ from $0.25$ to $0.1$ \emph{increased} TPF from $1.409$ to
$1.441$.

At convergence the objective decomposes as in Table~\ref{tab:decomp}: the
verifier term is $98.2\%$ of the total.

\begin{table}[!tbh]
\centering
\begin{tabular}{lrrrr}
\toprule
Term & Value & Weight & Contribution & Share \\
\midrule
$\mathcal{L}_{\mathrm{ver}}$ & 1.8480 & 1.00 & 1.8480 & 98.2\% \\
$\mathcal{L}_{\mathrm{EC}}$  & 0.0670 & 0.40 & 0.0268 & 1.4\% \\
$\mathcal{L}_{\mathrm{TD}}$  & 0.0250 & 0.20 & 0.0050 & 0.3\% \\
$\mathcal{L}_{\mathrm{end}}$ & 0.0109 & 0.25 & 0.0027 & 0.1\% \\
\bottomrule
\end{tabular}
\caption{Validation loss decomposition at the end of the main run.}
\label{tab:decomp}
\end{table}

\paragraph{What this does and does not establish.} The decomposition is a
snapshot at convergence and small residuals are what a satisfied constraint
looks like; it does not by itself prove that a term was unnecessary. Part of
the endpoint term's decay is moreover trivial: as $t \to 1$ the input $x_t$
approaches $x_1$ and predicting the endpoint degenerates towards copying the
input. Separating a learned denoiser from a learned copy requires the ablation
$w_{\mathrm{end}} = 0$, and separating the flow-map contribution from plain
verifier distillation requires $w_{\mathrm{ver}} = 0$ at matched budget.
Neither has been run.

\cm{Required before submission, in priority order.
\emph{Evaluation only, no retraining:}
(1) re-evaluate both existing checkpoints on the held-out benchmarks with a
sample size that supports a confidence interval --- without this the headline
number is inside baseline noise and the paper is rejectable on that alone;
(2) jump-count sweep $n \in \{1,2,4\}$ on the existing FlowDraft checkpoint ---
a stated deliverable, never run, and the only evidence that the two-time family
is used at inference;
(3) Orthrus reproduction validated against the authors' released weights, using
the harness already in the repository.
\emph{Requires training:}
(4) rerun the geometry pair of Table~\ref{tab:geometry} at the budget of the
main experiment --- it is the only controlled comparison in the paper and it
currently rests on 209 steps and two evaluation prompts;
(5) $w_{\mathrm{ver}} = 0$ at matched budget --- decides whether the paper has a
flow-map claim at all;
(6) $\lambda = 0$ at matched budget --- isolates self-distillation;
(7) $\ge 3$ seeds per arm at matched \texttt{anchors\_per\_sequence} and lr.
Without (1) no comparative claim survives review; without (2) and (4) the
abstract must not attribute the gain to the flow-map formulation.}

\subsection{Positional acceptance}

\cm{TODO: figure --- per-position acceptance within the block for both
drafters. Data is logged as \texttt{val/acceptance\_pos\_NN}; the baseline
profile decays $0.914, 0.652, 0.477, 0.366, 0.294, 0.240, 0.193$ across the
seven drafted positions. Comparing the two profiles shows whether the flow-map
drafter gains uniformly or specifically at the far positions, which is the
distinction between a better joint proposal and a better marginal one.}

% =========================================================================
\section{Limitations}
\label{sec:limitations}

We state the limitations of the present evidence explicitly.

\emph{Measurement.} Decode metrics are computed over a small set of validation
prompts. The run-to-run spread of the baseline under an identical configuration
($\pm 9\%$ TPF) is larger than the improvement we report over its best run
($+3.9\%$), so the comparison in Table~\ref{tab:main} should be read as
indicative rather than established. \cm{Remove or rewrite once item~(1) of the
required work is done.}

\emph{Attribution.} The self-distillation terms cannot be credited with the
gain on the present evidence (Section~\ref{sec:attribution}), and the ablations
that would isolate the flow-map contribution have not been run. We therefore
claim a drafter that performs at least as well as the masked-diffusion
baseline, and an analysis of which term drives it --- not that categorical flow
maps as such improve speculative decoding.

\emph{Inference regime.} All reported numbers use a one-jump schedule. The
two-time family $\pmap{s}{t}$ for $s>0$ is exercised during training but never
at inference, so the few-step regime that motivates the flow-map construction
is untested here. Under the time-sampling scheme of
Section~\ref{sec:objective} the operating corner $(0,1)$ is moreover a
measure-zero event, and $\mathcal{L}_{\mathrm{ver}}$ is the only term
supervising it directly; the family away from that corner receives no verifier
signal at all.

\emph{Scale.} A single backbone (Qwen3-1.7B), a single block size for the main
comparison, and a single seed per arm.

% =========================================================================
\section{Future Work}
\label{sec:future}

The limitation of Section~\ref{sec:limitations} that is structural rather than
budgetary is that the two-time family is trained but never composed. Every term
of Eq.~\eqref{eq:loss} evaluates a single jump; the only mechanism connecting
jumps is $\mathcal{L}_{\mathrm{EC}}$, and it connects them through a
stop-gradient target. A decode schedule with $n>1$ therefore executes a
composition that no term of the objective has ever seen end to end.

Two directions follow, in increasing order of departure from the present
construction.

\paragraph{Supervising the last-jump family.} Every schedule, whatever its
length, terminates on $\pmap{s}{1}$ at some intermediate noise level, and only
$s=0$ receives verifier signal today. Drawing $s$ across $[0,1)$ --- with
stratification, since the family is wide and the batch is small --- extends the
verifier term to that whole family at no additional forward cost, being the
same forward under a different time label. This is implemented and awaits
evaluation. Note that the same construction must not be extended across $t$:
because $\pmap{s}{t}$ is an endpoint predictor, a target that does not depend
on $t$ would drive the family towards $t$-independence and collapse exactly the
structure it is meant to train.

\paragraph{Differentiating through the rollout.} A stronger version replaces
the single-jump evaluation with an actual multi-jump rollout at training time:
sample a schedule length, roll the drafter forward through it, accumulate the
loss along the trajectory, and backpropagate through the chain of calls. This
trains the composition directly rather than hoping it follows from pairwise
consistency. The costs are real and should be stated: memory grows with
schedule length, each step consumes the previous step's learned output so
errors compound and the gradient inherits the conditioning of a recurrent
computation, and --- most importantly --- the stop-gradient that
$\mathcal{L}_{\mathrm{EC}}$ relies on is precisely what makes the
self-distillation bound available, so a fully differentiated rollout trades a
guarantee for a stronger signal. \cm{The structural analogy to volumetric
rendering, where stratified samples along a ray are composited and
differentiated through, is suggestive but incomplete: there the field is static
and the samples are independent, whereas here each sample is the previous
sample's output under the same trained map. Decide before submission whether to
draw this comparison at all.}

% =========================================================================
\section{Conclusion}
\label{sec:conclusion}

We replaced the masked-diffusion drafter of a lossless parallel decoder with a
categorical flow map trained in the geometry it decodes in, and obtained
$2.07$ accepted tokens per cycle against $2.00$ for the reproduced baseline, at
identical output. We also traced which part of the objective produces that
gain, and found that the verifier alignment term accounts for the bulk of it
while the self-distillation terms reach full weight after the gain is already
realised. We regard the second result as the more useful of the two: it
delimits what may presently be attributed to the flow-map formulation, and
identifies the ablations that would settle the question.

\cm{Rewrite once the ablations land. If $w_{\mathrm{ver}} = 0$ holds up, this
becomes a flow-map paper; if it does not, the honest framing is a simplex
parameterisation of the Orthrus drafter plus an attribution analysis, and the
title should change accordingly.}

% =========================================================================
\section{Contribution}
\label{sec:contribution}

\paragraph{Scope inherited from the project brief.} The research direction ---
using a categorical flow map as the drafter inside the Orthrus
lossless-verification scaffold, and designing a dual distillation objective
combining AR-teacher alignment with flow-map consistency --- was set by the
project brief, as were the four deliverables: reproduce Orthrus, implement the
flow-map drafter, develop the dual distillation objective, and evaluate the
three decoders on acceptance length, TPF and throughput with block-size and
jump-count ablations. The work below is the team's execution of that brief and
its extensions beyond it.

\paragraph{Ya.~Sergaev.} Built the codebase: the Orthrus adapter over a frozen
HuggingFace backbone, the diffusion head, the flow-time embedding, the simplex
preprocessor, both drafter variants (full-sequence and block-wise), the
speculative decoding loop with greedy and coupled-sampling verification, and
the evaluation harness together with the acceptance, TPF and losslessness
metrics. Designed the experimental matrix --- baseline, baseline in block
geometry, staged flow map, flow map in block geometry, and the two objective
ablations. Contributed, beyond the brief, the block-wise formulation in which
the drafter is trained in the exact geometry it decodes in
(Section~\ref{sec:geometry}), which is the configuration all reported results
use. Identified and diagnosed the failure of the initial objective, in which
the verifier term anchored the diagonal: established from training curves that
acceptance collapsed to zero while teacher agreement kept improving, and
proposed that the AR-teacher and self-distillation terms were mutually
incompatible as placed. Performed the attribution analysis of
Section~\ref{sec:attribution} and wrote the paper.

\paragraph{V.~Tekaev.} Reformulated the objective so that the diagonal is
anchored by the categorical endpoint likelihood rather than the AR teacher, and
relocated the verifier term to the jump the decode loop executes
(Section~\ref{sec:objective}) --- the change that made the drafter train.
Implemented inference-aligned training with multi-anchor blocks, packed
training with gradient checkpointing, multi-GPU training, and the numerical
stability and attention-backend fixes required at scale. Produced the campaign
configurations and ran all experiments reported in
Section~\ref{sec:results}, including the on-policy positional acceptance
metrics.

\paragraph{N.~Nikonov.} Built the validation harness for the Orthrus
reproduction, comparing our implementation against the authors' released
weights and code, which is what substantiates the baseline column of
Table~\ref{tab:main}.

\cm{Fourth team member: add a line here or remove them from the author list.
The call makes per-member contribution a review criterion, so an unexplained
name is worse than a short honest line.}

\cm{Cross-check this section against the repository history before submission:
several files were renamed mid-project, so \texttt{git log} on a current
filename attributes the file to whoever renamed it unless \texttt{-{}-follow}
is passed. See \texttt{PROVENANCE.md}.}

\section*{Acknowledgments}
\cm{Omit in the review version.}

\providecommand{\bysame}{\leavevmode\hbox to3em{\hrulefill}\thinspace}
\providecommand{\href}[2]{#2}
\begin{thebibliography}{10}

\bibitem{davis2026scaling}
\cm{TODO: fill from arXiv:2605.07820}, \emph{Scaling categorical flow maps},
  arXiv preprint arXiv:2605.07820, 2026.

\bibitem{nguyen2026orthrus}
Chien~Van Nguyen, Chaitra Hegde, Van~Cuong Pham, Ryan~A. Rossi, Franck
  Dernoncourt, and Thien~Huu Nguyen, \emph{Orthrus: Memory-efficient parallel
  token generation via dual-view diffusion}, arXiv preprint arXiv:2605.12825,
  2026.

\bibitem{roos2026cfm}
Daan Roos, Oscar Davis, Floor Eijkelboom, Michael Bronstein, Max Welling,
  İsmail İlkan Ceylan, Luca Ambrogioni, and Jan-Willem van~de Meent,
  \emph{Categorical flow maps}, Proceedings of the International Conference on
  Machine Learning (ICML), 2026, arXiv:2602.12233,
  \url{https://arxiv.org/abs/2602.12233}.

\end{thebibliography}

\vspace{1cm}
\selectlanguage{russian}
\noindent\textbf{Аноним}

\smallskip
\noindent\textbf{FlowDraft: драфтер на категориальных картах потока для
несмещённого параллельного декодирования}

\smallskip
\noindent\textbf{Ключевые слова:} спекулятивное декодирование, карты потока,
дискретная диффузия, эффективный инференс, языковые модели.

\smallskip
\noindent\cm{TODO: русская аннотация, перевод английской.}
\selectlanguage{english}

\end{document}
